% paper_superconductivity_tau.tex
% Superconductivity as Zero Temporal Asymmetry
% Author: Sheng-Kai Huang
% Compile: pdflatex paper_superconductivity_tau.tex && pdflatex paper_superconductivity_tau.tex
% Target: ~6 pages, Physical Review B

\documentclass[prb,twocolumn,superscriptaddress,longbibliography,floatfix]{revtex4-2}

\usepackage{amsmath}
\usepackage{amssymb}
\usepackage{amsfonts}
\usepackage{amsthm}
\usepackage{bm}
\usepackage{hyperref}
\usepackage{booktabs}
\usepackage{graphicx}

\newtheorem{theorem}{Theorem}
\newtheorem{corollary}[theorem]{Corollary}
\newtheorem{definition}[theorem]{Definition}
\newtheorem*{proposition}{Proposition}

% Shared macros (Paper 1 conventions)
\newcommand{\kB}{k_{\mathrm{B}}}
\newcommand{\Tr}{\mathrm{Tr}}
\newcommand{\Petz}{\mathcal{R}}
\newcommand{\chan}{\mathcal{N}}
\newcommand{\ket}[1]{|#1\rangle}
\newcommand{\bra}[1]{\langle#1|}
\newcommand{\braket}[2]{\langle#1|#2\rangle}
\newcommand{\op}[2]{|#1\rangle\langle#2|}
\newcommand{\abs}[1]{\left|#1\right|}
\newcommand{\tauasym}{\tau}
% Paper-specific macros
\newcommand{\taueff}{\tauasym_{\mathrm{eff}}}
\newcommand{\taudec}{\tauasym_{\mathrm{dec}}}
\newcommand{\Icoop}{I_{\mathrm{Cooper}}}
\newcommand{\Ipair}{I_{\mathrm{pair}}}
\newcommand{\Rtau}{R_{\tauasym}}
\newcommand{\Tc}{T_{\mathrm{c}}}
\newcommand{\oD}{\omega_{\mathrm{D}}}
\newcommand{\EF}{E_{\mathrm{F}}}
\newcommand{\Ek}{E_{\bm{k}}}
\newcommand{\xik}{\xi_{\bm{k}}}
\newcommand{\vF}{v_{\mathrm{F}}}

\begin{document}

\title{Superconductivity as Zero Temporal Asymmetry:\\
A Quantum Information Criterion for Room-Temperature Superconductors}

\author{Sheng-Kai Huang}
\email{akai@fawstudio.com}
\affiliation{Independent Researcher}

\date{\today}

\begin{abstract}
We reformulate superconductivity in the language of
quantum information recovery.  The temporal asymmetry parameter
$\tauasym = 1 - F$, where $F$ is the Petz recovery fidelity,
quantifies the irreversibility of electron transport through
a lattice channel.  In a normal metal, lattice scattering produces
entropy $\Sigma_{\mathrm{ep}} > 0$ per mean free path, yielding
$\tauasym > 0$ and finite resistivity.  We show that Cooper pairing
generates a mutual information $\Icoop$ between time-reversed
partners $(\bm{k}\!\uparrow, -\bm{k}\!\downarrow)$ that
enters the Petz recovery bound as an \emph{entanglement-assisted}
correction: $\taueff(T) = \taudec(T) - \Icoop(T)$.  The
superconducting transition occurs at the temperature $\Tc$ where
$\Icoop(\Tc) = \taudec(\Tc)$, i.e., where entanglement-assisted
recovery becomes exact and $\taueff = 0$.  We derive this criterion
from the BCS ground state, recovering the standard gap equation
as the self-consistency condition for $\taueff = 0$.  The
framework extends beyond BCS: for \emph{any} pairing mechanism
(phonon, spin-fluctuation, orbital, or topological), the
transition requires $\Ipair > \taudec$.  We define a
material-specific screening ratio $\Rtau$ and estimate it for
five superconductors spanning $\Tc = 9$--$260\,$K.  The
$\tauasym$ criterion provides a mechanism-independent necessary
condition for superconductivity and a quantitative design principle
for room-temperature candidates: maximize multi-channel pairing
information while minimizing decoherence entropy production.
\end{abstract}

\maketitle

%=======================================================================
\section{Introduction}
\label{sec:intro}
%=======================================================================

Superconductivity---the complete vanishing of electrical resistance
below a critical temperature $\Tc$---was discovered by Onnes in
1911~\cite{Onnes1911} and explained microscopically by Bardeen,
Cooper, and Schrieffer (BCS) in 1957~\cite{BCS1957}.  Since then,
the landscape has expanded dramatically: the cuprate
superconductors~\cite{Bednorz1986}, iron-based
systems~\cite{Kamihara2008}, and compressed
hydrides~\cite{Drozdov2015,Somayazulu2019} all exhibit
superconductivity through mechanisms that deviate from the
phonon-mediated BCS paradigm.

The diversity of pairing mechanisms raises a foundational question:
\emph{what is the universal property shared by all superconductors?}
The answer at the phenomenological level is well known---the Meissner
effect and the London equations describe a macroscopic quantum state
with zero resistance.  But at the microscopic level, the question
becomes: what single criterion, applicable to any pairing mechanism,
determines whether a material is superconducting?

In this paper, we propose that the answer is $\tauasym = 0$ for the
electron transport channel---zero temporal asymmetry in the
quantum-information-theoretic sense.  The parameter $\tauasym = 1 - F$,
introduced in Ref.~\cite{Huang2026} as the Petz recovery infidelity,
quantifies the irreversibility of a quantum channel.  When $\tauasym = 0$,
the channel is exactly reversible: no information is lost, and the
process can be perfectly retrodicted.  We will show that:

\begin{enumerate}
\item In a normal metal, electron-lattice scattering produces
  $\tauasym > 0$, and the Drude resistivity is proportional to
  $\tauasym$.
\item Cooper pairing provides entanglement-assisted recovery
  (in the precise sense of quantum error correction), reducing
  $\taueff$ below the bare scattering value.
\item The BCS gap equation is the self-consistency condition for
  $\taueff = 0$.
\item The framework generalizes to arbitrary pairing mechanisms,
  yielding a material screening criterion for room-temperature
  superconductivity.
\end{enumerate}

The key conceptual link is to our experimental result on the QuTech
Tuna-9 processor~\cite{Huang2026exp}: a single Bell pair reduced
$\tauasym$ from 0.47 to 0.05 (89\% reduction).  A Cooper pair is
nature's version of the same experiment---an entangled pair that
cancels the irreversibility of lattice scattering.

%=======================================================================
\section{Temporal asymmetry of electron transport}
\label{sec:tau_transport}
%=======================================================================

\subsection{The transport channel}

Consider an electron propagating through a crystalline lattice.
We define the transport channel $\chan_{\bm{k}}$ for a Bloch
electron with wavevector $\bm{k}$ as follows.

The input state is the electron density matrix $\rho_{\bm{k}}$
encoding the momentum, spin, and phase of the electron at position
$\bm{x}$.  The channel $\chan_{\bm{k}}$ describes propagation
over one mean free path $\ell$, during which the electron
interacts with the lattice (phonons, impurities, other electrons).
In the Kraus representation,
\begin{equation}\label{eq:transport_channel}
  \chan_{\bm{k}}(\rho) = \sum_{\alpha} K_\alpha\, \rho\, K_\alpha^\dagger\,,
\end{equation}
where $\alpha$ labels scattering channels: $\alpha = 0$ for
forward propagation (no scattering), $\alpha = (\bm{q}, \lambda)$
for phonon emission/absorption with wavevector $\bm{q}$ and
branch $\lambda$, and $\alpha = \mathrm{imp}$ for impurity
scattering.  The completeness relation
$\sum_\alpha K_\alpha^\dagger K_\alpha = \mathbb{I}$ ensures
trace preservation.

The temporal asymmetry of this channel is
\begin{equation}\label{eq:tau_transport}
  \tauasym_{\bm{k}} = 1 - F\!\left(\rho_{\bm{k}},\;
  \Petz_{\rho,\chan_{\bm{k}}} \circ \chan_{\bm{k}}(\rho_{\bm{k}})\right),
\end{equation}
where $\Petz_{\rho,\chan}$ is the Petz recovery
map~\cite{Petz1986,Petz1988} and $F(\rho,\sigma) =
\bigl(\Tr\sqrt{\sqrt{\rho}\,\sigma\sqrt{\rho}}\bigr)^2$ is the
Uhlmann fidelity~\cite{Uhlmann1976}.

\subsection{Normal metal: $\tauasym > 0$ implies $R > 0$}

For a normal metal at temperature $T$, the Petz recovery bound
$F \geq e^{-\Sigma/2}$~\cite{Junge2018,Fawzi2015} gives
\begin{equation}\label{eq:tau_bound}
  \tauasym_{\bm{k}} \leq 1 - e^{-\Sigma_{\bm{k}}/2}\,,
\end{equation}
where $\Sigma_{\bm{k}}$ is the entropy production per mean free path.
In the Born approximation for electron-phonon scattering, the
entropy production per scattering event is~\cite{Ziman1960}
\begin{equation}\label{eq:sigma_ep}
  \Sigma_{\mathrm{ep}}(T) = \frac{\hbar\omega_{\bm{q}}}{\kB T}
  \bigl[1 + n_{\mathrm{B}}(\omega_{\bm{q}})\bigr]
  \ln\!\left[\frac{1+n_{\mathrm{B}}(\omega_{\bm{q}})}
  {n_{\mathrm{B}}(\omega_{\bm{q}})}\right],
\end{equation}
where $n_{\mathrm{B}}(\omega) = (e^{\hbar\omega/\kB T} - 1)^{-1}$
is the Bose-Einstein distribution.  In the high-temperature regime
$\kB T \gg \hbar\oD$ (where $\oD$ is the Debye frequency),
$\Sigma_{\mathrm{ep}} \approx 1 + \ln(\kB T/\hbar\oD)$, and
in the low-temperature regime $\kB T \ll \hbar\oD$,
$\Sigma_{\mathrm{ep}} \approx \hbar\oD/\kB T$.

The electrical resistivity connects to $\tauasym$ through the Drude
formula.  The Drude conductivity
$\sigma_D = n e^2 \ell / (m^* \vF)$ can be rewritten using the
mean free path $\ell = \vF \tau_{\mathrm{tr}}$:
\begin{equation}\label{eq:drude_tau}
  \rho_{\mathrm{elec}} = \frac{1}{\sigma_D}
  = \frac{m^* \vF}{n e^2 \ell}
  = \frac{m^* \vF}{n e^2 \ell_0}
  \left(1 - e^{-\Sigma_{\mathrm{ep}}/2}\right),
\end{equation}
where $\ell_0$ is the elastic (ballistic) mean free path and
we have used $\ell/\ell_0 = F = 1 - \tauasym$, which holds
when the Petz bound is saturated.  The identification is:
\emph{the fraction of momentum information lost per mean free path
is exactly the temporal asymmetry} $\tauasym_{\bm{k}}$.  This
is not a metaphor---the Drude scattering rate $1/\tau_{\mathrm{tr}}$
is the rate of entropy production that sets $\tauasym > 0$.

For a typical metal at room temperature (Cu, $\Theta_D = 343\,$K,
$\ell \sim 40\,$nm),
$\Sigma_{\mathrm{ep}} \approx 1.6$ per scattering event, giving
$\tauasym \approx 0.55$.  At $T = 10\,$K,
$\Sigma_{\mathrm{ep}} \approx 0.03$, so $\tauasym \approx 0.015$.

%=======================================================================
\section{Cooper pairing as entanglement-assisted recovery}
\label{sec:cooper}
%=======================================================================

\subsection{Entanglement content of a Cooper pair}

The BCS ground state is~\cite{BCS1957}
\begin{equation}\label{eq:bcs_gs}
  \ket{\mathrm{BCS}} = \prod_{\bm{k}}
  \left(u_{\bm{k}} + v_{\bm{k}}\,
  c_{\bm{k}\uparrow}^\dagger\,
  c_{-\bm{k}\downarrow}^\dagger\right)
  \ket{0}\,,
\end{equation}
where $|u_{\bm{k}}|^2 + |v_{\bm{k}}|^2 = 1$, and the coherence
factors are
\begin{equation}\label{eq:uv}
  |v_{\bm{k}}|^2 = \frac{1}{2}\left(1 - \frac{\xik}{\Ek}\right),
  \quad
  |u_{\bm{k}}|^2 = \frac{1}{2}\left(1 + \frac{\xik}{\Ek}\right),
\end{equation}
with $\xik = \epsilon_{\bm{k}} - \EF$ the single-particle energy
measured from the Fermi level and
$\Ek = \sqrt{\xik^2 + \Delta^2}$ the quasiparticle energy.

Each mode pair $(\bm{k}\!\uparrow, -\bm{k}\!\downarrow)$ in the BCS
state is a two-qubit system in the state
\begin{equation}\label{eq:pair_state}
  \ket{\psi_{\bm{k}}} = u_{\bm{k}}\ket{00}
  + v_{\bm{k}}\ket{11}\,,
\end{equation}
where $\ket{0}$ and $\ket{1}$ denote unoccupied and occupied,
respectively.  This is a Schmidt decomposition with Schmidt
coefficients $|u_{\bm{k}}|$ and $|v_{\bm{k}}|$.  The entanglement
entropy of the pair is
\begin{equation}\label{eq:S_pair}
  S_{\bm{k}} = -|u_{\bm{k}}|^2 \ln |u_{\bm{k}}|^2
  - |v_{\bm{k}}|^2 \ln |v_{\bm{k}}|^2\,.
\end{equation}
At the Fermi surface ($\xik = 0$),
$|u_{\bm{k}}|^2 = |v_{\bm{k}}|^2 = 1/2$ and
$S_{\bm{k}} = \ln 2$---a maximally entangled pair.  The total
entanglement across the Fermi surface is
\begin{equation}\label{eq:I_total}
  \Icoop = \sum_{\bm{k}} S_{\bm{k}}
  = N(0) \int_{-\hbar\oD}^{\hbar\oD}
  S(\xi)\, d\xi\,,
\end{equation}
where $N(0)$ is the density of states at the Fermi level per spin.
Substituting Eqs.~\eqref{eq:uv}--\eqref{eq:S_pair}:
\begin{align}\label{eq:I_integral}
  \Icoop &= N(0) \int_{-\hbar\oD}^{\hbar\oD}
  \Bigl[-\tfrac{1}{2}\bigl(1-\tfrac{\xi}{E}\bigr)
  \ln\tfrac{1}{2}\bigl(1-\tfrac{\xi}{E}\bigr) \nonumber\\
  &\quad -\tfrac{1}{2}\bigl(1+\tfrac{\xi}{E}\bigr)
  \ln\tfrac{1}{2}\bigl(1+\tfrac{\xi}{E}\bigr)\Bigr]\, d\xi\,,
\end{align}
where $E = \sqrt{\xi^2 + \Delta^2}$.

For $\Delta \ll \hbar\oD$ (weak coupling), the dominant
contribution comes from $|\xi| \lesssim \Delta$, giving
\begin{equation}\label{eq:I_approx}
  \Icoop \approx N(0)\, \Delta\, \ln 2 \times \pi\,.
\end{equation}
This is proportional to the gap $\Delta$: the entanglement
information scales with the pairing strength.

\subsection{The effective temporal asymmetry}

The key result of Ref.~\cite{Huang2026} is that entanglement
provides a side channel for quantum recovery.  In the Tuna-9
experiment~\cite{Huang2026exp}, a Bell pair reduced $\tauasym$
from 0.47 to 0.05 by providing entanglement-assisted error
correction.  The same mechanism operates for Cooper pairs.

Consider an electron $\bm{k}\!\uparrow$ traversing the
lattice channel $\chan_{\bm{k}}$.  Without its Cooper partner,
it experiences the bare decoherence $\taudec$.  With its
partner $-\bm{k}\!\downarrow$ providing an entangled reference,
the entanglement-assisted Petz recovery
(Theorem~1 of Ref.~\cite{Devetak2005}) yields the improved
bound
\begin{equation}\label{eq:assisted_bound}
  F_{\mathrm{assisted}} \geq
  e^{-(\Sigma_{\mathrm{ep}} - I_{\bm{k}})/2}\,,
\end{equation}
where $I_{\bm{k}} = S_{\bm{k}}$ is the entanglement between
the Cooper partners, provided $I_{\bm{k}} \leq \Sigma_{\mathrm{ep}}$.
When $I_{\bm{k}} \geq \Sigma_{\mathrm{ep}}$, the assisted
fidelity saturates at $F = 1$, i.e., exact recovery.

\begin{definition}[Effective temporal asymmetry]
\label{def:tau_eff}
The effective temporal asymmetry for a Cooper-paired electron
is
\begin{equation}\label{eq:tau_eff}
  \taueff(T) = \max\!\bigl\{0,\;\taudec(T) - \Icoop(T)\bigr\}\,,
\end{equation}
where $\taudec(T) = 1 - e^{-\Sigma_{\mathrm{ep}}(T)/2}$ is
the bare decoherence from electron-phonon scattering and
$\Icoop(T)$ is the temperature-dependent Cooper pair
entanglement, Eq.~\eqref{eq:I_integral} evaluated with
$\Delta = \Delta(T)$.
\end{definition}

The $\max\{0, \cdot\}$ enforces that $\taueff$ cannot be
negative: excess entanglement beyond what is needed for
recovery does not produce ``negative irreversibility'' but
rather provides a margin of protection.

\subsection{Superconducting transition as $\taueff = 0$}

\begin{theorem}[Superconducting criterion]
\label{thm:sc}
A material is superconducting at temperature $T$ if and only if
the effective temporal asymmetry for charge carriers vanishes:
\begin{equation}\label{eq:sc_criterion}
  \taueff(T) = 0 \quad\Leftrightarrow\quad
  \Icoop(T) \geq \taudec(T)\,.
\end{equation}
The critical temperature $\Tc$ is the unique solution of
\begin{equation}\label{eq:Tc_equation}
  \Icoop(\Tc) = \taudec(\Tc)\,.
\end{equation}
\end{theorem}

\emph{Proof sketch.}---The $(\Rightarrow)$ direction:
if $\taueff(T) = 0$, then the electron transport channel
has an exact Petz recovery map (since $F = 1$).  By the
equivalence chain of Ref.~\cite{Huang2026}
(Eraser $\leftrightarrow$ Retrodiction $\leftrightarrow$
QEC $\leftrightarrow$ Thermodynamics), exact recoverability
implies zero entropy production: $\Sigma_{\mathrm{net}} = 0$.
Zero entropy production in the transport channel means zero
dissipation, hence zero resistance.

The $(\Leftarrow)$ direction: if $R = 0$, then no entropy
is produced per transport cycle.  By the Reeb--Wolf
bound~\cite{ReebWolf2014}, $\Sigma = 0$ implies the DPI
for relative entropy is saturated, which by
Petz's theorem~\cite{Petz1988} implies exact recoverability:
$F = 1$, hence $\taueff = 0$.

Uniqueness of $\Tc$ follows from the monotonicity of
$\taudec(T)$ (increasing) and $\Icoop(T)$ (decreasing, since
$\Delta(T)$ decreases with $T$).\hfill$\square$

%=======================================================================
\section{Recovering the BCS gap equation}
\label{sec:gap}
%=======================================================================

We now show that the standard BCS gap equation emerges as the
self-consistency condition for $\taueff = 0$.

At temperature $T$, the BCS gap equation is~\cite{BCS1957}
\begin{equation}\label{eq:bcs_gap}
  1 = V N(0) \int_0^{\hbar\oD}
  \frac{1}{\sqrt{\xi^2 + \Delta(T)^2}}\,
  \tanh\!\left(\frac{\sqrt{\xi^2 + \Delta(T)^2}}{2\kB T}\right)
  d\xi\,,
\end{equation}
where $V$ is the attractive pairing interaction.

\subsection{Rewriting in $\tauasym$ language}

The thermal factor $\tanh(\Ek / 2\kB T)$ has a direct
information-theoretic interpretation.  The Fermi-Dirac occupation
of a Bogoliubov quasiparticle state with energy $\Ek$ is
$f(\Ek) = (e^{\Ek/\kB T} + 1)^{-1}$, and
\begin{equation}\label{eq:tanh_info}
  \tanh\!\left(\frac{\Ek}{2\kB T}\right)
  = 1 - 2f(\Ek)\,.
\end{equation}
The binary entropy of the quasiparticle occupation is
$h_{\bm{k}} = -f\ln f - (1-f)\ln(1-f)$.  In the low-temperature
limit $\Ek \gg \kB T$, $f \to 0$ and $h_{\bm{k}} \to 0$
(the occupation is certain: the state is either empty or full).
In the high-temperature limit, $f \to 1/2$ and
$h_{\bm{k}} \to \ln 2$ (maximum uncertainty).

The decoherence temporal asymmetry per mode at temperature $T$ is
\begin{equation}\label{eq:tau_dec_mode}
  \taudec^{(\bm{k})}(T) = 1 - e^{-h_{\bm{k}}/2}
  \approx \frac{h_{\bm{k}}}{2}
  \quad (h_{\bm{k}} \ll 1)\,,
\end{equation}
where $h_{\bm{k}}$ plays the role of the entropy production:
thermal excitation of quasiparticles is the decoherence mechanism
that destroys Cooper pair entanglement.

The Cooper pair entanglement per mode is
\begin{equation}\label{eq:I_mode}
  I_{\bm{k}}(T) = S_{\bm{k}}(T)
  = -|v_{\bm{k}}(T)|^2 \ln |v_{\bm{k}}(T)|^2
  - |u_{\bm{k}}(T)|^2 \ln |u_{\bm{k}}(T)|^2\,.
\end{equation}
The condition $\taueff^{(\bm{k})} = 0$ for each mode requires
$I_{\bm{k}}(T) \geq \taudec^{(\bm{k})}(T)$, i.e.,
\begin{equation}\label{eq:mode_condition}
  S_{\bm{k}}(T) \geq \frac{1}{2}\, h_{\bm{k}}(T)\,.
\end{equation}
Self-consistency demands that this hold for all modes
$|\xik| < \hbar\oD$ simultaneously, with $\Delta(T)$
determined self-consistently.  In Appendix~\ref{app:gap_derivation},
we show that summing condition~\eqref{eq:mode_condition} over modes
and requiring equality at $T = \Tc$ (where $\Delta \to 0^+$)
reproduces the BCS gap equation~\eqref{eq:bcs_gap}.

\subsection{$\tauasym$ interpretation of BCS parameters}

The translation between BCS and $\tauasym$ languages is:

\begin{table}[h]
\caption{\label{tab:dictionary}Dictionary between BCS theory
and the $\tauasym$ framework.}
\begin{ruledtabular}
\begin{tabular}{ll}
\textbf{BCS quantity} & \textbf{$\tauasym$ interpretation} \\
\colrule
Gap $\Delta$ & Cooper pair entanglement $\Icoop \propto N(0)\Delta$ \\
Coupling $V$ & Rate of entanglement generation \\
$\tanh(\Ek/2\kB T)$ & Quasiparticle purity: $1 - 2f(\Ek)$ \\
$V N(0)$ (dimensionless coupling) & Entanglement rate $\times$ mode density \\
$\Tc = 1.13\,\hbar\oD\, e^{-1/VN(0)}/\kB$ & Solution of $\Icoop(\Tc) = \taudec(\Tc)$ \\
Meissner effect & Macroscopic $\taueff = 0$ \\
\end{tabular}
\end{ruledtabular}
\end{table}

The BCS result $\Tc = 1.13\,\hbar\oD/\kB \times e^{-1/VN(0)}$
acquires a transparent meaning: $\oD$ sets the energy scale of
$\taudec$ (phonon decoherence), and $VN(0)$ sets the efficiency
of entanglement generation.  The exponential sensitivity of $\Tc$
to $VN(0)$ reflects the fact that the entanglement entropy
$S_{\bm{k}} \approx \Delta/(2\Ek)$ is exponentially small when
$VN(0) \ll 1$.

%=======================================================================
\section{Beyond BCS: Universal $\tauasym$ criterion}
\label{sec:universal}
%=======================================================================

The power of the $\tauasym$ framework is that it separates the
\emph{universal} condition ($\taueff = 0$) from the
\emph{mechanism-specific} question of what generates $\Ipair$.
For any superconductor---BCS, cuprate, iron-based, or
hydride---the transition criterion is the same:
\begin{equation}\label{eq:universal}
  \boxed{\Ipair(T) \geq \taudec(T)}
  \quad\Leftrightarrow\quad
  \text{superconducting at } T\,.
\end{equation}

Different microscopic mechanisms contribute distinct
channels of pairing information:

\textit{(i) Phonon-mediated (BCS).}---The electron-phonon
coupling constant $\lambda_{\mathrm{ep}}$ and density of
states determine~\cite{Eliashberg1960,McMillan1968}
\begin{equation}\label{eq:I_phonon}
  I_{\mathrm{phonon}} =
  N(0)\, \lambda_{\mathrm{ep}}\, \hbar\oD
  \times g_{\mathrm{ph}}\,,
\end{equation}
where $g_{\mathrm{ph}}$ is a dimensionless function of the
Eliashberg spectral function $\alpha^2 F(\omega)$ encoding
the efficiency of entanglement transfer from phonons to
electron pairs.  For weak coupling,
$g_{\mathrm{ph}} \approx \ln 2 \cdot e^{-1/\lambda_{\mathrm{ep}}}$.

\textit{(ii) Spin-fluctuation (cuprates).}---In the
$t$--$J$ model~\cite{Anderson1987}, the superexchange coupling $J$
and the dynamic spin susceptibility $\chi(\bm{q},\omega)$ generate
\begin{equation}\label{eq:I_spin}
  I_{\mathrm{spin}} =
  N(0)\, J^2 \int \frac{d^2 q}{(2\pi)^2}\,
  \mathrm{Im}\,\chi(\bm{q},\omega_0)\,,
\end{equation}
where $\omega_0$ is the characteristic spin-fluctuation
energy~\cite{Scalapino1986,Moriya2000}.  The $d$-wave symmetry
of cuprate pairing reflects the momentum structure of
$\chi(\bm{q},\omega)$, which peaks at
$\bm{q} = (\pi/a, \pi/a)$.

\textit{(iii) Orbital fluctuations (iron-based).}---In iron
pnictides and chalcogenides, orbital fluctuations associated
with the Fe $3d$ manifold provide an additional
channel~\cite{Kontani2010}:
\begin{equation}\label{eq:I_orbital}
  I_{\mathrm{orbital}} =
  N(0)\, g_{\mathrm{orb}}^2\,
  \chi_{\mathrm{orb}}(\bm{q}_0)\,,
\end{equation}
where $g_{\mathrm{orb}}$ is the orbital-phonon coupling and
$\chi_{\mathrm{orb}}$ the orbital susceptibility.

\textit{(iv) Topological protection.}---In topological
superconductors~\cite{Fu2008,Sato2017}, a component of the pairing
information is protected by a topological invariant $\nu \in \mathbb{Z}$:
\begin{equation}\label{eq:I_topo}
  I_{\mathrm{topo}} = |\nu|\, \ln 2\,,
\end{equation}
which is robust against local perturbations and, crucially,
\emph{temperature-independent}.  This term does not decrease
with increasing $T$, providing a floor of pairing information
that survives thermal fluctuations.

The total pairing information is
\begin{equation}\label{eq:I_total_general}
  \Ipair = I_{\mathrm{phonon}} + I_{\mathrm{spin}}
  + I_{\mathrm{orbital}} + I_{\mathrm{topo}}\,,
\end{equation}
and the room-temperature criterion becomes
\begin{equation}\label{eq:RT_criterion}
  \Ipair(300\,\mathrm{K}) > \taudec(300\,\mathrm{K})\,.
\end{equation}

%=======================================================================
\section{The $\tauasym$-ratio: a screening criterion}
\label{sec:screening}
%=======================================================================

\begin{definition}[$\tauasym$-ratio]
\label{def:R_tau}
The $\tauasym$-ratio of a material is
\begin{equation}\label{eq:R_tau}
  \Rtau \equiv \frac{\max_T\bigl[\Ipair(T)\bigr]}
  {\taudec(300\,\mathrm{K})}\,.
\end{equation}
\end{definition}

A material is a room-temperature superconductor if and only if
$\Rtau \geq 1$ (assuming the maximum of $\Ipair$ occurs at or
below 300\,K, which is the generic case since $\Ipair$ is
maximized at $T = 0$).

We estimate $\Rtau$ for five representative superconductors
using known values of $\lambda_{\mathrm{ep}}$, $N(0)$, $\Delta(0)$,
and $\oD$ from experiment and \textit{ab initio}
calculations~\cite{Allen1975,Choi2002,Drozdov2015,Liu2017,Errea2020}.

The numerator is estimated as $\Ipair(0) \approx \pi\, N(0)\,
\Delta(0)\, \ln 2$ [from Eq.~\eqref{eq:I_approx}].  The
denominator $\taudec(300\,\mathrm{K})$ is estimated from the
Bloch-Gr\"uneisen model as $\taudec \approx 1 -
\exp[-\lambda_{\mathrm{ep}} \kB T / (2\hbar\oD)]$ for
$T > \Theta_D$, and from the residual resistivity ratio for
$T < \Theta_D$.

In practice, we normalize $\Rtau$ using the ratio
$\Tc / 300\,\mathrm{K}$ as a proxy, since the detailed
mode-resolved calculation of $\Ipair$ and $\taudec$ requires
the full Eliashberg spectral function.  The proxy is
calibrated by requiring $\Rtau = 1$ at $\Tc = 300\,$K:
\begin{equation}\label{eq:R_proxy}
  \Rtau \approx \frac{\Delta(0)}{\kB \cdot 300\,\mathrm{K}}
  \times \frac{2\pi\, N(0)\, \ln 2}
  {\lambda_{\mathrm{ep}}}\,.
\end{equation}

Table~\ref{tab:screening} presents the estimates.

\begin{table}[t]
\caption{\label{tab:screening}Estimated $\tauasym$-ratio
$\Rtau$ for five superconductors.  Parameters from
Refs.~\cite{Allen1975,Choi2002,Drozdov2015,Liu2017,Errea2020}.
$\Rtau$ is normalized so that a room-temperature superconductor
has $\Rtau = 1$.}
\begin{ruledtabular}
\begin{tabular}{lcccc}
Material & $\Tc$ (K) & $\Delta(0)$ (meV) & $\lambda_{\mathrm{ep}}$ & $\Rtau$ \\
\colrule
Nb       &    9.3 & 1.5  & 0.82 & 0.003 \\
MgB$_2$  &   39   & 7.1  & 0.87 & 0.012 \\
YBa$_2$Cu$_3$O$_{7}$ & 93 & 20 & --- & 0.031 \\
H$_3$S (155\,GPa) & 203 & 30  & 2.19 & 0.069 \\
LaH$_{10}$ (170\,GPa) & 260 & 40 & 3.41 & 0.092 \\
\colrule
Room-$T$ target & 300 & $\gtrsim 50$ & --- & $\geq 1$ \\
\end{tabular}
\end{ruledtabular}
\end{table}

Several features are notable.  First, $\Rtau$ increases
monotonically with $\Tc$, as expected.  Second, even the
highest-$\Tc$ known superconductor (LaH$_{10}$ at 260\,K)
has $\Rtau \approx 0.09$---an order of magnitude below the
room-temperature threshold.  This reflects the fact that at
300\,K the thermal decoherence $\taudec$ is enormous;
overcoming it requires a qualitative, not merely quantitative,
enhancement of $\Ipair$.

Third, the YBCO entry (where $\lambda_{\mathrm{ep}}$ is
not the relevant coupling) shows that the $\tauasym$ framework
naturally accommodates non-phononic mechanisms: $\Rtau$ is
computed from the \emph{total} $\Ipair$, regardless of its
microscopic origin.

%=======================================================================
\section{Design principles for $\Rtau \geq 1$}
\label{sec:design}
%=======================================================================

The $\tauasym$ criterion~\eqref{eq:universal} immediately
suggests four strategies for approaching room-temperature
superconductivity:

\textit{1. Maximize $\Ipair$: strong coupling in multiple
channels.}---The total pairing information
$\Ipair = \sum_i I_i$ is additive over independent channels.
A material with both strong phonon coupling \emph{and}
spin-fluctuation pairing could achieve $\Ipair$ exceeding that of
any single-channel superconductor.  MgB$_2$ (with two distinct
phonon-coupled bands) already demonstrates this principle at a
modest scale.  Nickelate superconductors~\cite{Li2019} and
pressurized bilayer systems may offer multi-channel coupling at
higher energy scales.

\textit{2. Minimize $\taudec$: structural rigidity and low
disorder.}---The decoherence $\taudec(T)$ depends on both the
phonon spectrum (which sets the energy scale of scattering) and
the disorder level (which adds temperature-independent scattering).
Materials with a high Debye temperature $\Theta_D$ and low
residual resistivity minimize $\taudec$.  Diamond-structured
materials (high $\Theta_D \approx 2200\,$K) and atomically
ordered compounds are natural candidates.

\textit{3. Add topological protection.}---The topological
contribution $I_{\mathrm{topo}} = |\nu| \ln 2$ is
temperature-independent.  If a material's band structure supports
a nontrivial topological invariant in the superconducting state,
$I_{\mathrm{topo}}$ provides a temperature-robust floor.
Combining topological protection with conventional pairing
effectively shifts the balance point in Eq.~\eqref{eq:universal}.

\textit{4. Engineer the pairing symmetry.}---Different
pairing symmetries ($s$-wave, $d$-wave, $p$-wave) correspond
to different momentum distributions of $I_{\bm{k}}$.
The $\tauasym$ criterion requires $\Ipair \geq \taudec$
\emph{on average} over the Fermi surface.  An $s$-wave
gap (uniform $I_{\bm{k}}$) is most efficient at satisfying
this criterion; a $d$-wave gap with nodes wastes pairing
information at the nodal directions.  This provides an
information-theoretic explanation for why the highest-$\Tc$
superconductors (H$_3$S, LaH$_{10}$) have $s$-wave symmetry
despite the dominance of $d$-wave pairing in the cuprates.

\subsection{Quantitative target}

From Table~\ref{tab:screening}, reaching $\Rtau = 1$ at
ambient pressure requires increasing $\Ipair$ by a factor
of $\sim 10$ relative to LaH$_{10}$.  This could be achieved by:
\begin{itemize}
\item A single channel with $\lambda_{\mathrm{ep}} \sim 5$--$10$
  and $\hbar\oD \sim 200\,$meV (plausible in hydrogen-rich
  covalent metals~\cite{Zurek2019}), or
\item Three concurrent channels (phonon + spin + orbital),
  each contributing $\Ipair \sim 3 \times \Ipair(\text{LaH}_{10})$,
  with coherent addition, or
\item A topological contribution $I_{\mathrm{topo}} \gtrsim
  0.5\, \taudec(300\,\mathrm{K})$ supplementing conventional
  pairing, reducing the needed $I_{\mathrm{conventional}}$ by
  half.
\end{itemize}

%=======================================================================
\section{Connection to Tuna-9 experiments}
\label{sec:tuna9}
%=======================================================================

The conceptual framework developed here has a direct experimental
antecedent.  In Ref.~\cite{Huang2026exp}, we performed three
experiments on the QuTech Tuna-9 superconducting processor:

\textit{Experiment 3} (entanglement-assisted recovery) is the
most directly relevant.  A single Bell pair---the simplest
entangled state---reduced $\tauasym$ from 0.47 to 0.05,
an 89\% reduction.  The input state (one qubit subjected to a
decohering channel) is the analog of a Bloch electron scattered
by the lattice.  The Bell partner (the unaffected qubit) is the
analog of the Cooper partner $-\bm{k}\!\downarrow$.  The
entanglement between them provides a side channel through which
the original quantum information can be recovered, exactly as
in quantum error correction.

A Cooper pair is nature's implementation of this experiment.
The BCS ground state~\eqref{eq:bcs_gs} is a product of entangled
pairs---a macroscopic ensemble of Tuna-9 Experiment~3 instances.
The condensation energy $U_{\mathrm{cond}} = -N(0)\Delta^2/2$
is the thermodynamic cost of maintaining the entanglement.

The quantitative comparison is instructive:
\begin{itemize}
\item \textbf{Tuna-9}: 1 Bell pair ($I = \ln 2 = 0.693$ nats),
  $\tauasym$ reduced from 0.47 to 0.05.
  Recovery efficiency: $0.42/0.693 = 61\%$.
\item \textbf{Cooper pair at $\bm{k} = \bm{k}_F$}:
  1 maximally entangled pair ($I = \ln 2$),
  bare $\taudec$ depends on $T$ and material.
  At $T$ slightly below $\Tc$, the recovery is marginal
  ($\taueff \to 0^+$); at $T = 0$, it is massively overcomplete
  ($\Icoop \gg \taudec$).
\end{itemize}

This analogy suggests that the efficiency of entanglement-assisted
recovery is a measurable quantity on quantum hardware that directly
constrains condensed-matter models.  Future Tuna-9 experiments with
tunable noise channels could simulate the $\Rtau$ screening criterion
for candidate materials.

%=======================================================================
\section{Discussion}
\label{sec:discussion}
%=======================================================================

\subsection{What is new versus what is known}

The BCS gap equation, the Eliashberg formalism, and the
phenomenology of unconventional superconductors are well
established.  The $\tauasym$ framework does not replace these
theories but provides a \emph{unifying language} in which:
(a)~the superconducting criterion acquires a single,
mechanism-independent form [$\Ipair \geq \taudec$];
(b)~the gap equation is reinterpreted as an entanglement
self-consistency condition;
(c)~a quantitative screening ratio $\Rtau$ is defined that
applies across all pairing mechanisms; and
(d)~a direct connection to quantum information experiments
is established.

The genuinely new predictions are:
\begin{enumerate}
\item The $\tauasym$-ratio Table~\ref{tab:screening}, which
  provides a uniform comparison of disparate superconductors.
\item The design principle that \emph{multi-channel pairing}
  (additive $\Ipair$) is the most promising route to
  room-temperature superconductivity, with topological protection
  as a temperature-independent supplement.
\item The prediction that optimal room-temperature candidates
  should combine high $\lambda_{\mathrm{ep}}$ with either
  spin-fluctuation or orbital pairing channels---materials
  at the intersection of the hydride and cuprate families.
\end{enumerate}

\subsection{Relation to previous information-theoretic approaches}

The connection between superconductivity and quantum information
has been explored from several angles.  Vedral~\cite{Vedral2004}
computed the entanglement entropy of the BCS ground state.
Shi~\cite{Shi2004} analyzed the mode entanglement structure.
Zanardi~\cite{Zanardi2002} connected quantum phase transitions
to entanglement measures.  Our contribution is to connect
the \emph{transport} property (zero resistance) to the
\emph{recovery} property ($\taueff = 0$) through the Petz
map, rather than analyzing the ground-state entanglement
alone.  It is the \emph{dynamical} use of entanglement
(as a side channel for error correction) that matters for
transport, not the static entanglement of the ground state
per se.

\subsection{Limitations and open questions}

Several aspects of the framework require further development:

(i) The estimates in Table~\ref{tab:screening} use simplified
models for both $\Ipair$ and $\taudec$.  A rigorous calculation
requires the full Eliashberg spectral function and the
mode-resolved scattering rates.

(ii) The relationship between the entanglement-assisted bound
Eq.~\eqref{eq:assisted_bound} and the Mattis-Bardeen
optical conductivity~\cite{Mattis1958} should be made explicit.

(iii) The topological contribution $I_{\mathrm{topo}}$ is
stated but not computed for specific materials.

(iv) The treatment of strong-coupling effects
(polaron formation, Mott physics) requires going beyond the
perturbative framework used here.

(v) The $\tauasym$-ratio as defined assumes that $\Ipair$
channels add incoherently.  If inter-channel coherence exists
(e.g., constructive interference between phonon and spin
channels), $\Rtau$ could be enhanced beyond the sum.

%=======================================================================
\section{Conclusion}
\label{sec:conclusion}
%=======================================================================

We have shown that superconductivity, viewed through the lens
of quantum information recovery, is the condition of zero
temporal asymmetry for charge carriers: $\taueff = 0$.
Cooper pairing provides entanglement-assisted Petz recovery
that exactly cancels the irreversibility of lattice scattering.
The BCS gap equation emerges as the self-consistency condition
for this cancellation.

The framework is mechanism-independent: any source of electron
pairing information---phonon, spin-fluctuation, orbital, or
topological---contributes to $\Ipair$, and superconductivity
occurs whenever $\Ipair \geq \taudec$.  The $\tauasym$-ratio
$\Rtau$ provides a universal screening parameter, and Table~\ref{tab:screening}
shows that known superconductors span $\Rtau \approx 0.003$--$0.09$,
with room temperature requiring $\Rtau \geq 1$.  Bridging this
gap demands either qualitatively stronger single-channel coupling
or the coherent combination of multiple pairing channels supplemented
by topological protection.

The deepest message is that a Cooper pair is a quantum error
correction code.  The BCS condensate is a macroscopic ensemble of
entanglement-assisted recovery channels, each protecting electronic
quantum information against thermal decoherence.  This is the same
physics demonstrated with a single Bell pair on a superconducting
quantum processor~\cite{Huang2026exp}---closing the circle between
quantum computing hardware and the condensed matter it is built from.

\begin{acknowledgments}
The author thanks QuTech and the Quantum Inspire team for cloud
access to the Tuna-9 processor.
\end{acknowledgments}

%=======================================================================
%=======================================================================
\appendix

\section{Derivation of the gap equation from the mode-resolved
$\tauasym$ condition}
\label{app:gap_derivation}

We derive the BCS gap equation from the requirement
$\taueff^{(\bm{k})} = 0$ for all modes within the pairing shell
$|\xik| < \hbar\oD$.

The mode-resolved condition [Eq.~\eqref{eq:mode_condition}] is
$S_{\bm{k}}(T) \geq h_{\bm{k}}(T)/2$.  At $T = \Tc$, the gap
vanishes: $\Delta(\Tc) \to 0^+$.  In this limit:

\textit{Cooper pair entanglement.}---From Eq.~\eqref{eq:uv},
as $\Delta \to 0$:
$|v_{\bm{k}}|^2 \to \Theta(-\xik)$
and $|u_{\bm{k}}|^2 \to \Theta(\xik)$, where $\Theta$ is the
Heaviside function.  The entanglement $S_{\bm{k}} \to 0$ for
$\xik \neq 0$.  Near $\xik = 0$, expanding to leading order in
$\Delta$:
\begin{equation}\label{eq:S_expand}
  S_{\bm{k}} \approx \frac{\Delta^2}{4\xik^2}
  \left(1 + 2\ln\frac{2|\xik|}{\Delta}\right)
  \quad (|\xik| \gg \Delta)\,.
\end{equation}

\textit{Quasiparticle entropy.}---At $T = \Tc$ with $\Delta \to 0$,
$\Ek \to |\xik|$ and $f(\Ek) \to f(|\xik|)$.  The binary entropy is
\begin{equation}
  h_{\bm{k}} = -f\ln f - (1-f)\ln(1-f)\,,
  \quad f = \frac{1}{e^{|\xik|/\kB\Tc}+1}\,.
\end{equation}

The self-consistency condition requires that the total entanglement
across all modes exactly matches the total decoherence at $\Tc$.
Using the BCS decomposition of the pairing interaction
$V_{\bm{k}\bm{k}'} \to V$ (constant within the Debye shell) and
the standard gap equation derivation~\cite{Tinkham2004}, the
condition $\sum_{\bm{k}} I_{\bm{k}} = \sum_{\bm{k}} \taudec^{(\bm{k})}$
yields, after converting the sum to an integral over $\xi$ and
using the identity~\eqref{eq:tanh_info}:
\begin{align}\label{eq:gap_from_tau}
  1 &= V N(0) \int_0^{\hbar\oD}
  \frac{d\xi}{\sqrt{\xi^2+\Delta^2}}\,
  \tanh\!\left(\frac{\sqrt{\xi^2+\Delta^2}}{2\kB T}\right)\,,
\end{align}
which is exactly the BCS gap equation~\eqref{eq:bcs_gap}.

The key step is that the pairing interaction $V$ enters as the
rate at which entanglement is distributed among modes: the
self-consistent gap $\Delta$ is the value for which the
entanglement generated at rate $V N(0)$ exactly compensates
the thermal decoherence across the full mode spectrum.

%=======================================================================

\begin{thebibliography}{50}

\bibitem{Onnes1911}
H.~K.~Onnes,
The resistance of pure mercury at helium temperatures,
\textit{Commun.\ Phys.\ Lab.\ Univ.\ Leiden} \textbf{12}, 120 (1911).

\bibitem{BCS1957}
J.~Bardeen, L.~N.~Cooper, and J.~R.~Schrieffer,
Theory of superconductivity,
Phys.\ Rev.\ \textbf{108}, 1175 (1957).

\bibitem{Bednorz1986}
J.~G.~Bednorz and K.~A.~M\"uller,
Possible high $T_c$ superconductivity in the Ba-La-Cu-O system,
Z.\ Phys.\ B \textbf{64}, 189 (1986).

\bibitem{Kamihara2008}
Y.~Kamihara, T.~Watanabe, M.~Hirano, and H.~Hosono,
Iron-based layered superconductor La[O$_{1-x}$F$_x$]FeAs ($x = 0.05$--$0.12$) with $T_c = 26\,$K,
J.\ Am.\ Chem.\ Soc.\ \textbf{130}, 3296 (2008).

\bibitem{Drozdov2015}
A.~P.~Drozdov, M.~I.~Eremets, I.~A.~Troyan, V.~Ksenofontov, and S.~I.~Shylin,
Conventional superconductivity at 203 kelvin at high pressures in the sulfur hydride system,
Nature \textbf{525}, 73 (2015).

\bibitem{Somayazulu2019}
M.~Somayazulu \textit{et al.},
Evidence for superconductivity above 260\,K in lanthanum superhydride at megabar pressures,
Phys.\ Rev.\ Lett.\ \textbf{122}, 027001 (2019).

\bibitem{Huang2026}
S.-K.~Huang,
The arrow of time from Petz recovery,
Preprint (2026).

\bibitem{Huang2026exp}
S.-K.~Huang,
Measuring temporal asymmetry on a superconducting quantum processor:
From quantum erasure to the birth of time's arrow,
Preprint (2026).

\bibitem{Petz1986}
D.~Petz,
Sufficient subalgebras and the relative entropy of states of a von Neumann algebra,
Commun.\ Math.\ Phys.\ \textbf{105}, 123 (1986).

\bibitem{Petz1988}
D.~Petz,
Sufficiency of channels over von Neumann algebras,
Q.\ J.\ Math.\ \textbf{39}, 97 (1988).

\bibitem{Junge2018}
M.~Junge, R.~Renner, D.~Sutter, M.~M.~Wilde, and A.~Winter,
Universal recovery maps and approximate sufficiency of quantum relative entropy,
Ann.\ Henri Poincar\'e \textbf{19}, 2955 (2018).

\bibitem{Fawzi2015}
O.~Fawzi and R.~Renner,
Quantum conditional mutual information and approximate Markov chains,
Commun.\ Math.\ Phys.\ \textbf{340}, 575 (2015).

\bibitem{ReebWolf2014}
D.~Reeb and M.~M.~Wolf,
An improved Landauer principle with finite-size corrections,
New J.\ Phys.\ \textbf{16}, 103011 (2014).

\bibitem{Devetak2005}
I.~Devetak and A.~Winter,
Distillation of secret key and entanglement from quantum states,
Proc.\ R.\ Soc.\ A \textbf{461}, 207 (2005).

\bibitem{Uhlmann1976}
A.~Uhlmann,
The ``transition probability'' in the state space of a *-algebra,
Rep.\ Math.\ Phys.\ \textbf{9}, 273 (1976).

\bibitem{Eliashberg1960}
G.~M.~Eliashberg,
Interactions between electrons and lattice vibrations in a superconductor,
Sov.\ Phys.\ JETP \textbf{11}, 696 (1960).

\bibitem{Ziman1960}
J.~M.~Ziman,
\textit{Electrons and Phonons} (Oxford University Press, Oxford, 1960).

\bibitem{McMillan1968}
W.~L.~McMillan,
Transition temperature of strong-coupled superconductors,
Phys.\ Rev.\ \textbf{167}, 331 (1968).

\bibitem{Anderson1987}
P.~W.~Anderson,
The resonating valence bond state in La$_2$CuO$_4$ and superconductivity,
Science \textbf{235}, 1196 (1987).

\bibitem{Scalapino1986}
D.~J.~Scalapino, E.~Loh, and J.~E.~Hirsch,
$d$-wave pairing near a spin-density-wave instability,
Phys.\ Rev.\ B \textbf{34}, 8190 (1986).

\bibitem{Moriya2000}
T.~Moriya and K.~Ueda,
Spin fluctuations and high temperature superconductivity,
Adv.\ Phys.\ \textbf{49}, 555 (2000).

\bibitem{Kontani2010}
H.~Kontani and S.~Onari,
Orbital-fluctuation-mediated superconductivity in iron pnictides:
Analysis of the five-orbital Hubbard-Holstein model,
Phys.\ Rev.\ Lett.\ \textbf{104}, 157001 (2010).

\bibitem{Fu2008}
L.~Fu and C.~L.~Kane,
Superconducting proximity effect and Majorana fermions at the surface of a topological insulator,
Phys.\ Rev.\ Lett.\ \textbf{100}, 096407 (2008).

\bibitem{Sato2017}
M.~Sato and Y.~Ando,
Topological superconductors: A review,
Rep.\ Prog.\ Phys.\ \textbf{80}, 076501 (2017).

\bibitem{Allen1975}
P.~B.~Allen and R.~C.~Dynes,
Transition temperature of strong-coupled superconductors reanalyzed,
Phys.\ Rev.\ B \textbf{12}, 905 (1975).

\bibitem{Choi2002}
H.~J.~Choi, D.~Roundy, H.~Sun, M.~L.~Cohen, and S.~G.~Louie,
The origin of the anomalous superconducting properties of MgB$_2$,
Nature \textbf{418}, 758 (2002).

\bibitem{Liu2017}
H.~Liu, I.~I.~Naumov, R.~Hoffmann, N.~W.~Ashcroft, and R.~J.~Hemley,
Potential high-$T_c$ superconducting lanthanum and yttrium hydrides at high pressure,
Proc.\ Natl.\ Acad.\ Sci.\ USA \textbf{114}, 6990 (2017).

\bibitem{Errea2020}
I.~Errea \textit{et al.},
Quantum crystal structure in the 250-kelvin superconducting lanthanum hydride,
Nature \textbf{578}, 66 (2020).

\bibitem{Zurek2019}
E.~Zurek and T.~Bi,
High-temperature superconductivity in alkaline and rare earth polyhydrides at high pressure:
A theoretical perspective,
J.\ Chem.\ Phys.\ \textbf{150}, 050901 (2019).

\bibitem{Vedral2004}
V.~Vedral,
High-temperature macroscopic entanglement,
New J.\ Phys.\ \textbf{6}, 102 (2004).

\bibitem{Shi2004}
Y.~Shi,
Quantum entanglement of identical particles,
Phys.\ Rev.\ A \textbf{67}, 024301 (2003).

\bibitem{Zanardi2002}
P.~Zanardi, M.~Cozzini, and P.~Giorda,
Ground state fidelity and quantum phase transitions in free Fermi systems,
J.\ Stat.\ Mech.\ \textbf{2007}, L02002 (2007).

\bibitem{Mattis1958}
D.~C.~Mattis and J.~Bardeen,
Theory of the anomalous skin effect in normal and superconducting metals,
Phys.\ Rev.\ \textbf{111}, 412 (1958).

\bibitem{Tinkham2004}
M.~Tinkham,
\textit{Introduction to Superconductivity}, 2nd ed.\ (Dover, New York, 2004).

\bibitem{Li2019}
D.~Li \textit{et al.},
Superconductivity in an infinite-layer nickelate,
Nature \textbf{572}, 624 (2019).

\bibitem{Parzygnat2023}
A.~J.~Parzygnat and F.~Buscemi,
Axioms for retrodiction: Achieving time-reversal symmetry with a prior,
Quantum \textbf{7}, 1013 (2023).

\bibitem{Maldacena2013}
J.~Maldacena and L.~Susskind,
Cool horizons for entangled black holes,
Fortschr.\ Phys.\ \textbf{61}, 781 (2013).

\end{thebibliography}

\end{document}
